Considereu el sistema d'equacions lineals següent:
\[
\left.\begin{aligned}
y-z&=p+3\\
p^2x-z&=5\\
x-y&=3
\end{aligned}\right\}
\]

\begin{apartats}

\apartat{1,25}
Discutiu el sistema per als diferents valors del paràmetre $p$.

\begin{solucio}
La matriu de coeficients i l'ampliada, $A$ i $A'$ respectivament, són les següents:
\[
\left(\begin{array}{ccc|c}0&1&-1&p+3\\p^2&0&-1&5\\1&-1&0&3\end{array}\right).
\]
En primer lloc, calculem el determinant de $A$:
\[
|A|=\begin{vmatrix}0&1&-1\\p^2&0&-1\\1&-1&0\end{vmatrix}=0+p^2-1=p^2-1 .
\]
Igualant-lo a zero, en resulten els valors $p=1$ i $p=-1$. Per tant, distingim els casos següents:\\
Si $p\neq1,\,-1$, el determinant de $A$ és diferent de zero i, per tant, $\operatorname{rang}(A)=\operatorname{rang}(A')=3$,
que és igual al nombre d'incògnites. Per tant, es tracta d'un sistema compatible determinat.\\
Si $p=1$, el determinant de la matriu de coeficients és zero; com que tenim el menor
$\begin{vmatrix}0&1\\1&0\end{vmatrix}=-1\neq0$, aleshores $\operatorname{rang}(A)=2$. Però la matriu ampliada té
$\operatorname{rang}(A')=3$ gràcies, per exemple, al menor
\[
\begin{vmatrix}0&1&4\\1&0&5\\1&-1&3\end{vmatrix}=5-4-3=-2\neq0 .
\]
Per tant, el sistema és incompatible.\\
Si $p=-1$, tenim el sistema $y-z=2$, $x-z=5$, $x-y=3$. Eliminant la segona equació (per ser suma de la
primera i la tercera), i com que la matriu
\[
\begin{pmatrix}0&1&-1&2\\1&-1&0&3\end{pmatrix}
\]
té rang 2, el sistema és compatible indeterminat amb un grau de llibertat.

\textit{Pauta oficial:} 0,25 pel càlcul del determinant; 0,25 per determinar els valors crítics de $p$, i 0,25
per la discussió de cadascun dels tres casos. Hi ha altres maneres de discutir i resoldre el sistema
d'equacions: compteu-les bé en la mesura que ho facin correctament i de manera justificada.
\end{solucio}

\apartat{0,5}
Resoleu el sistema per al cas $p=-1$.

\begin{solucio}
Per al cas $p=-1$, es tracta del sistema $y-z=2$, $x-y=3$, que és compatible indeterminat. Fent $y=\lambda$,
s'obté $x=\lambda+3$ i $z=\lambda-2$. Per tant, la solució del sistema és, en aquest cas,
$(x,y,z)=(\lambda+3,\,\lambda,\,\lambda-2)$, amb $\lambda$ real.

\textit{Pauta oficial:} 0,5 pel càlcul correcte del conjunt de solucions (evidentment, poden triar qualsevol
de les variables com a lliure).
\end{solucio}

\apartat{0,75}
Per al cas $p=-1$, hi ha alguna solució que compleixi, a més, $xy=10$? En cas afirmatiu, indiqueu quantes
n'hi ha i trobeu-les totes.

\begin{solucio}
Per al cas $p=-1$, el sistema té les infinites solucions obtingudes a l'apartat anterior. Imposant, a més,
que $xy=10$, resulta $(\lambda+3)\lambda=10$, és a dir, $\lambda^2+3\lambda-10=0$; d'aquí obtenim
\[
\lambda=\frac{-3\pm\sqrt{9-4(-10)}}{2}=\frac{-3\pm7}{2}=-5,\ 2 .
\]
Per tant, el sistema té exactament dues solucions que compleixen $xy=10$, que són les corresponents a
$\lambda=-5$ i $\lambda=2$. Es tracta de $(x,y,z)=(-2,-5,-7)$ i $(x,y,z)=(5,2,0)$, respectivament.

\textit{Pauta oficial:} 0,25 pel plantejament de l'equació quadràtica; 0,25 per la seva resolució, i 0,25 pel
càlcul dels dos punts resultants.
\end{solucio}

\end{apartats}
